%! Setting Up a Test for the Difference of Two Population Proportions — Unit 6, Lesson 6.10 — Warm-Up
\documentclass[10pt]{article}
\usepackage{apstats-article}
\usepackage{apstats-boxes}
\newcommand{\TallMath}[1]{$\displaystyle #1\rule[-1.4em]{0pt}{3.2em}$}

\begin{document}

\pageheader{Unit 6, Lesson 6.10}{Warm-Up}
\namedateperiod

\vspace{0.15in}

\textbf{Spiral: Setting Up a One-Sample Test.}
A researcher claims that more than 40\% of adults in a city use a ride-share app weekly.
In a random sample of 250 adults, 115 reported doing so.

\begin{enumerate}
  \item State $H_0$ and $H_a$ in words and symbols.
        \writelines{2}

  \item Which proportion do you use for the Large Counts condition in a one-sample test ---
        $\hat p$ or $p_0$? \blank{2cm}

  \item Check Large Counts: $n \cdot p_0 = $ \blank{2cm}, $\quad n(1-p_0) = $ \blank{2cm}. \quad Both $\ge 10$? \blank{1.5cm}

  \item \textbf{Bridge.} Now suppose we have \textit{two} groups and want to test whether
        their proportions are equal. Under $H_0: p_1 = p_2$, is it still correct to use
        $\hat p_1$ or $\hat p_2$ separately for Large Counts? Why or why not?
        \writelines{2}
\end{enumerate}

\end{document}
